Análisis
La distribución del saldo total administrado por las agencias presenta diferencias moderadas entre las cinco sucursales analizadas. La agencia de
Petén
concentra el mayor volumen de saldos, con aproximadamente
Q5.34 millones
, mientras que la agencia ubicada en
Zona 10
registra el menor saldo agregado, con alrededor de
Q4.71 millones
. La diferencia entre ambas asciende a aproximadamente
Q500 mil
, lo que representa una variación relativamente reducida considerando el volumen total administrado por cada sucursal.
En términos generales, no se observa una concentración significativa de los saldos en una única agencia. Los resultados reflejan una distribución relativamente equilibrada entre las sucursales, lo que sugiere que la cartera de clientes y los recursos administrados no dependen de forma predominante de una sola ubicación geográfica.
Un aspecto que merece un análisis más detallado es que las agencias ubicadas fuera del área metropolitana presentan niveles de saldos comparables e incluso superiores a la sucursal de Zona 10. Este comportamiento podría estar asociado a diversos factores, como la composición de la cartera de clientes, la presencia de actividades económicas relevantes en dichas regiones, el nivel de bancarización o estrategias comerciales específicas implementadas por la institución. No obstante, la información disponible no permite establecer una relación causal, por lo que sería necesario incorporar variables adicionales para explicar este patrón con mayor precisión.